\chapter{Preliminares matemáticos}
\label{cap:prelim}

Este capítulo es la caja de herramientas que el resto del libro usa sin
volver a justificar. Cubre tres temas: (i) grafos dirigidos tipados, la
estructura sobre la que vive sotFS; (ii) funciones de hash
criptográfico, pieza central del direccionamiento por contenido; y
(iii) pushouts y reescritura DPO, el formalismo en el que se enuncian
las operaciones POSIX.

Si el lector ya viene con experiencia en estos temas, puede saltearse
este capítulo y volver a él como referencia. El estilo es deliberadamente
operativo: definiciones precisas, ejemplos chicos, sin esoterismo
categórico más allá de lo necesario.

\section{Grafos dirigidos tipados}

\begin{defn}[Grafo dirigido]
\label{def:grafo-dirigido}
Un \emph{grafo dirigido} es una tupla $G = (V, E, \mathrm{src},
\mathrm{tgt})$ donde $V$ y $E$ son conjuntos finitos y
$\mathrm{src}, \mathrm{tgt}: E \to V$ son funciones que asignan a cada
arista su origen y destino.
\end{defn}

\begin{defn}[Morfismo de grafos]
\label{def:morfismo-grafos}
Un \emph{morfismo de grafos} $f: G \to H$ es un par de funciones
$f_V: V_G \to V_H$ y $f_E: E_G \to E_H$ tales que $f_V \circ
\mathrm{src}_G = \mathrm{src}_H \circ f_E$ y análogamente para
$\mathrm{tgt}$. Los morfismos preservan la estructura de adyacencia.
\end{defn}

\begin{defn}[Grafo tipado]
\label{def:grafo-tipado}
Un \emph{grafo tipado} sobre una signatura $(T_V, T_E, \sigma)$ es un
grafo dirigido más dos funciones $\tau_V: V \to T_V$ y $\tau_E: E \to
T_E$ que respetan $\sigma$: para toda $e \in E$, $\sigma(\tau_E(e)) =
(\tau_V(\mathrm{src}(e)), \tau_V(\mathrm{tgt}(e)))$. La signatura
$\sigma: T_E \to T_V \times T_V$ fija qué pares de tipos puede unir cada
tipo de arista.
\end{defn}

La categoría de los grafos tipados sobre una signatura fija se denota
$\TGraph$. Sus objetos son grafos tipados; sus morfismos son los
morfismos de grafos que además preservan los tipos
($\tau_V \circ f_V = \tau_V'$ y análogo para $\tau_E$).

\begin{observacion}[$\TGraph$ es completa y co-completa]
\label{obs:tgraph-cocompleta}
Para los efectos del libro alcanza con un hecho: $\TGraph$ admite
todos los pushouts (constructibles componente a componente: la
construcción en $\Set$, una vez en $V$ y otra vez en $E$, respeta los
tipos por construcción). Esto es lo que permite definir reescritura DPO
sobre grafos tipados sin sutilezas adicionales.
\end{observacion}

\section{Funciones de hash criptográfico}

\begin{defn}[Función de hash criptográfico]
\label{def:hash}
Una función $h: \{0,1\}^* \to \{0,1\}^n$ es una \emph{función de hash
criptográfico} si satisface las siguientes propiedades:
\begin{description}
  \item[Resistencia a pre-imagen.] Dado $y \in \{0,1\}^n$, encontrar
    $x$ tal que $h(x) = y$ requiere esfuerzo $\Theta(2^n)$.
  \item[Resistencia a segunda pre-imagen.] Dado $x_1$, encontrar
    $x_2 \neq x_1$ con $h(x_1) = h(x_2)$ requiere esfuerzo
    $\Theta(2^n)$.
  \item[Resistencia a colisión.] Encontrar cualquier par $x_1 \neq
    x_2$ con $h(x_1) = h(x_2)$ requiere esfuerzo $\Theta(2^{n/2})$
    (cota del paradoja del cumpleaños).
\end{description}
\end{defn}

\paragraph{La elección de sotFS: BLAKE3.} sotFS usa \blake{}
\cite{BLAKE3} con $n = 256$ bits. La probabilidad de colisión accidental
entre dos contenidos distintos es $\sim 2^{-128}$, valor menor que el
de un error cósmico no detectado en una RAM ECC. \blake{} además es
muy rápido (~6\,GiB/s en una CPU moderna single-thread) y paraleliza
naturalmente sobre la entrada vía un Merkle tree interno: importante
para el direccionamiento por contenido del cap.~\ref{cap:hashing}.

\section{Pushouts}

\subsection*{En $\Set$}

\begin{defn}[Pushout en $\Set$]
\label{def:pushout-set-prelim}
Dadas dos funciones $f: A \to B$ y $g: A \to C$ con dominio común $A$,
su \emph{pushout} es un objeto $D$ junto con dos funciones $h: B \to
D$ y $k: C \to D$ tales que $h \circ f = k \circ g$ y, para cualquier
otro $D'$ con flechas $h', k'$ que satisfagan la misma ecuación, existe
una única flecha $u: D \to D'$ con $u \circ h = h'$ y $u \circ k = k'$.
\end{defn}

\begin{figure}[h]
\centering
\begin{tikzpicture}[node distance=15mm and 22mm]
  \node (A) {$A$};
  \node[right=of A] (B) {$B$};
  \node[below=of A] (C) {$C$};
  \node[below=of B] (D) {$D$};
  \node[below right=of D] (Dp) {$D'$};
  \draw[arr] (A) -- (B) node[edgelbl, midway] {$f$};
  \draw[arr] (A) -- (C) node[edgelbl, midway] {$g$};
  \draw[arr] (B) -- (D) node[edgelbl, midway] {$h$};
  \draw[arr] (C) -- (D) node[edgelbl, midway] {$k$};
  \draw[arr, dashed] (D) -- (Dp) node[edgelbl, midway] {$u$};
  \draw[arr, bend left=25] (B) to (Dp);
  \draw[arr, bend right=25] (C) to (Dp);
\end{tikzpicture}
\caption{Diagrama universal del pushout. $D$ es el ``más pequeño''
objeto que cierra el cuadrado conmutativo: cualquier otro $D'$ que lo
cierre factoriza por $D$ via $u$.}
\label{fig:pushout-prelim}
\end{figure}

\begin{teorema}[Construcción del pushout en $\Set$]
\label{teo:pushout-set-construccion}
$D$ existe siempre y es único salvo isomorfismo. Concretamente,
$D = (B \sqcup C)/\sim$ donde $\sim$ es la relación de equivalencia
generada por $f(a) \sim g(a)$ para todo $a \in A$.
\end{teorema}

\begin{proof}[Esquema]
La existencia se verifica construyendo $D$ como en el enunciado y
chequeando la propiedad universal directamente. La unicidad sale del
hecho de que cualquier otro $D'$ con la misma propiedad universal
admite una flecha $u: D \to D'$ y una $u': D' \to D$, ambas
factorizando las inclusiones, y la unicidad de $u, u'$ obliga
$u' \circ u = \mathrm{id}_D$.
\end{proof}

\subsection*{En $\Graph$ y $\TGraph$}

La construcción del pushout en $\Graph$ es la misma idea, pero
duplicada: una vez en el conjunto de vértices y otra en el de aristas.
Se puede hacer porque las funciones $\mathrm{src}, \mathrm{tgt}: E \to
V$ son funciones de conjuntos, y los morfismos de grafos las preservan.
En $\TGraph$ además los tipos se preservan por construcción si la
signatura $\sigma$ es respetada por los morfismos de partida.

\section{Reescritura DPO (\emph{double-pushout})}

La reescritura DPO es la herramienta para describir transformaciones
locales de un grafo de manera \emph{algebraica}: una regla especifica
qué borrar, qué dejar y qué insertar; un \emph{paso DPO} aplica la
regla a un grafo \emph{host} preservando el resto.

\begin{defn}[Producción / regla DPO]
\label{def:dpo-prelim}
Una \emph{producción DPO} es un span en la categoría $\TGraph$:
\[
  \rho \;=\; \dpo{L}{K}{R}
\]
donde $K \to L$ y $K \to R$ son inclusiones (morfismos inyectivos).
$L$ es el patrón a remover; $R$ el patrón a insertar; $K$ el
\emph{glue} que sobrevive.
\end{defn}

\begin{defn}[Paso DPO]
\label{def:dpo-step-prelim}
Dado un grafo host $G$ y un \emph{matching} $m: L \to G$, un paso DPO
$G \xRightarrow{\rho, m} H$ existe si y sólo si:
\begin{enumerate}
  \item Existe un grafo $D$ y un morfismo $D \to G$ tal que el cuadrado
    izquierdo (con $L \to G$, $K \to L$ y $K \to D$) sea un pushout.
    Este $D$ es el \emph{pushout complement} y existe si y sólo si se
    cumple la \emph{gluing condition} (def.~\ref{def:gluing-prelim}).
  \item El cuadrado derecho (con $R$, $K \to R$, $K \to D$) tiene un
    pushout, que es $H$.
\end{enumerate}
\end{defn}

\begin{defn}[Gluing condition]
\label{def:gluing-prelim}
La \emph{gluing condition} para $\rho = \dpo{L}{K}{R}$ y matching $m: L
\to G$ exige dos cosas:
\begin{description}
  \item[\textsf{Identification.}] Si dos elementos distintos de $L$
    (nodos o aristas) se mapean al mismo elemento de $G$ vía $m$,
    deben ambos pertenecer a la imagen de $K \to L$.
  \item[\textsf{Dangling.}] Para todo nodo $v \in m(L) \setminus
    m(K)$ (un nodo que la regla quiere borrar), no debe haber aristas
    en $G$ incidentes a $v$ que no provengan de $m(L)$.
\end{description}
La condición Dangling es la responsable de los errores POSIX que
veremos en el cap.~\ref{cap:dpo}: ``no podés borrar un directorio que
todavía tiene entradas''.
\end{defn}

\begin{teorema}[Existencia y unicidad del paso DPO]
\label{teo:dpo-existencia}
Dada una producción $\rho$ y un matching $m$ que satisface la gluing
condition, existen $D$ y $H$ únicos salvo isomorfismo, y son los
construidos vía pushout complement y pushout final.
\end{teorema}

La prueba completa vive en el clásico \cite{Ehrig2006}; es una
aplicación directa de las propiedades universales del pushout.

\section{Por qué DPO y no SPO o ARS}

Existen otros formalismos para reescritura de grafos: SPO
(\emph{single-pushout}, \cite{Corradini2006}), term rewriting con
\emph{attribute systems}, sistemas de reescritura abstracta (ARS).
\emph{¿Por qué DPO en sotFS?}

\begin{itemize}
  \item DPO requiere la gluing condition explícita; un paso ``ilegal''
    \emph{no existe}, no es ``hace algo raro''. En SPO, el paso siempre
    existe pero puede dejar el grafo en un estado degenerado (aristas
    colgando que se reinterpretan). DPO es \emph{conservador} en el
    sentido de Hippocratic: si no se puede hacer bien, no se hace.
  \item DPO es la formulación más cercana a las pre-condiciones POSIX:
    \texttt{EEXIST}, \texttt{ENOTEMPTY} y \texttt{ENOTDIR} son
    instancias literales de gluing conditions o NACs.
  \item Las pruebas formales en Coq se simplifican porque cada paso
    requiere un certificado (la gluing condition explícita), lo que
    hace los teoremas de preservación de invariantes más directos.
\end{itemize}

\section{Notación que se usa todo el libro}

\begin{itemize}
  \item $G, H$: grafos tipados (instancias de $\TGraph$). $G$ típicamente
    es un host; $H$ un resultado tras un paso.
  \item $L, K, R$: los tres componentes de una producción DPO.
  \item $m: L \to G$ un matching.
  \item $\hashof{x}$: el hash \blake{} de $x$, 256 bits.
  \item $\tau_V, \tau_E$: las funciones de tipado.
  \item $\NInode, \NDir, \NCap, \NTxn, \NVer, \NBlk$: los seis tipos de nodo
    de sotFS (cap.~\ref{cap:typegraph}).
  \item $\econtains, \egrants, \edelegates, \ederivedfrom,
    \esupersedes, \epointsto, \ehasxattr$: los siete tipos de arista.
\end{itemize}

\section*{Resumen del capítulo}

\begin{itemize}
  \item Un grafo tipado es un grafo dirigido con tipos asignados a
    nodos y aristas, sujetos a una signatura.
  \item Una función de hash criptográfico tiene tres propiedades de
    resistencia; \blake{} con 256 bits es la elección de sotFS.
  \item Un pushout en $\Set$ es el cociente de la unión disjunta por la
    relación inducida por las flechas; en $\TGraph$ se construye
    componente a componente.
  \item Una producción DPO es un span $\dpo{L}{K}{R}$; un paso requiere
    matching y gluing condition; produce un grafo $H$ único salvo
    isomorfismo.
  \item DPO se elige sobre SPO porque la gluing condition explícita
    cuadra con las pre-condiciones POSIX y simplifica las pruebas.
\end{itemize}

\section*{Ejercicios}

\begin{ejercicio}[\dificultad{$\star$}]
\label{ej:prelim-1}
Defina ``morfismo de grafos'' en una línea.
\end{ejercicio}

\begin{ejercicio}[\dificultad{$\star$}]
\label{ej:prelim-2}
¿Cuál es la cota del paradoja del cumpleaños para una función de hash
de 256 bits? Exprese el resultado como una potencia de 2.
\end{ejercicio}

\begin{ejercicio}[\dificultad{$\star\star$}]
\label{ej:prelim-3}
Calcule el pushout en $\Set$ de $f: \{a,b\} \to \{1,2,3\}$ con $f(a) = 1$,
$f(b) = 2$, y $g: \{a,b\} \to \{x,y,z\}$ con $g(a) = x$, $g(b) = y$.
Liste explícitamente los elementos de $D$ y las funciones $h, k$.
\end{ejercicio}

\begin{ejercicio}[\dificultad{$\star\star$}]
\label{ej:prelim-4}
Dé un ejemplo de producción DPO sobre $\Graph$ (sin tipar) que añada
una arista entre dos nodos existentes. Especifique $L$, $K$, $R$ con
diagramas o conjuntos.
\end{ejercicio}

\begin{ejercicio}[\dificultad{$\star\star$}]
\label{ej:prelim-5}
Construya un caso donde la gluing condition Dangling falle: un host
$G$ con tres nodos y dos aristas, un matching $m$ que selecciona uno
de los nodos para borrarlo, y la arista incidente sobreviviente que
hace fallar la condición.
\end{ejercicio}

\begin{ejercicio}[\dificultad{$\star\star$}]
\label{ej:prelim-6}
Dé un ejemplo en sotFS donde la condición Identification de la gluing
condition es relevante. Pista: ¿qué pasa si una regla DPO mapea dos
\NDir{} distintos al mismo \NDir{} del host?
\end{ejercicio}

\begin{ejercicio}[\dificultad{$\star\star$}]
\label{ej:prelim-7}
Compare DPO y SPO en términos de qué pasa con un \opunlink{} de un
inode hard-linked múltiples veces.
\end{ejercicio}

\begin{ejercicio}[\dificultad{$\star\star\star$}]
\label{ej:prelim-8}
Demuestre que el pushout complement, cuando existe, es único salvo
isomorfismo. Pista: aplicar la propiedad universal del pushout
izquierdo dos veces.
\end{ejercicio}

\begin{ejercicio}[\dificultad{$\star\star\star$}]
\label{ej:prelim-9}
Lea el cap.~3 de \cite{Ehrig2006}. Encuentre dos teoremas (con número
y página) que justifiquen por qué $\TGraph$ es una categoría adecuada
para reescritura DPO.
\end{ejercicio}
